\documentclass[11pt,a4paper]{article}
\usepackage[T1]{fontenc}
\usepackage[utf8]{inputenc}
\usepackage{lmodern,microtype,amsmath,amssymb,amsthm,mathtools,booktabs,array,xcolor}
\usepackage[margin=24mm]{geometry}
\usepackage{enumitem}
\usepackage[hypertexnames=false]{hyperref}
\definecolor{ink}{HTML}{183C54}
\hypersetup{colorlinks=true,linkcolor=ink,citecolor=ink,urlcolor=ink,
pdftitle={Critical Free-Poisson Densities and Cubic Supercongruences},
pdfauthor={Research manuscript prepared for Keio with ChatGPT}}
\newtheorem{theorem}{Theorem}[section]
\newtheorem{lemma}[theorem]{Lemma}
\newtheorem{proposition}[theorem]{Proposition}
\newtheorem{corollary}[theorem]{Corollary}
\newtheorem{problem}[theorem]{Auxiliary problem}
\newcommand{\vp}{v_p}
\numberwithin{equation}{section}
\setlength{\emergencystretch}{3em}
\setlength{\parskip}{3pt}
\title{Critical Free-Poisson Densities\\and Cubic Supercongruences\\[4pt]
\large Complete arguments for named conjectures, with a priority audit}
\author{Research manuscript prepared for Keio with ChatGPT}
\date{15 November 2026}
\begin{document}
\maketitle
\begin{abstract}
We give complete arguments concerning two previously stated conjectures.
For the rate-one free compound Poisson law $\mu=\pi\boxtimes\nu$ with
$\nu(\{0\})=0$, we prove that its density is essentially unbounded in
every neighborhood of zero. This establishes the discontinuity assertion
in Hasebe--Sakuma's Conjecture 3.7. Its stronger assertion of divergence
to infinity is false: an explicit positive atomic jump distribution gives
both arbitrarily large and arbitrarily small positive density values
approaching zero. The construction uses a quantitative comparison between
the mass of small jumps and the error from the larger jumps.
Independently, we prove Bala's supercongruence for the original sequence
A176285, for every prime $p\ge5$ and all positive indices and levels.
A homogeneous rational weight transfers a Pfaff--Saalsch\"utz sum through
the primitive and divisible index sectors. We also formulate and prove
the corresponding general weight-transfer theorem. A short appendix
gives an exact negative Hankel minor for Zhu's unweighted Jacobi--Stirling
conjecture. These mathematical claims are accompanied by complete proofs;
worldwide novelty and fulfillment of every historical-priority condition
are not certified by the bounded literature review.
\end{abstract}

\section{Targets, results, and the role of the corpus}
The distinction between a proof and a priority claim is essential here.
The arguments below use established theorems, not conjectural statements
from the supplied corpus. They are research manuscripts for expert
scrutiny, not independently refereed or formally verified publications.

\begin{center}\small
\begin{tabular}{@{}p{.32\textwidth}p{.19\textwidth}p{.41\textwidth}@{}}
\toprule
Original target & Source date & Outcome proved here\\
\midrule
Hasebe--Sakuma, Conjecture 3.7 & Preprint 2015; revised 2016 & Discontinuity holds for every permitted $\nu$; the stronger limit fails even through positive-density points.\\[4pt]
Bala, A176285 & 11 July 2024 & The complete congruence modulo $p^{3r}$ holds for every $p\ge5$ and $n,r\ge1$.\\[4pt]
Zhu, unweighted Jacobi--Stirling row conjecture & Attributed to 2018; printed in 2024 & A leading $3\times3$ Hankel minor has coefficient $-4$.\\
\bottomrule
\end{tabular}
\end{center}

The repository snapshot was commit \texttt{38348353416f1f7f}, consulted
on 15 November 2026 \cite{VMRC}. Its probability manuscripts [122] and
[127] suggested separating boundary behavior of a transform from global
shape claims. The congruence manuscripts [133] and [134] suggested
separating primitive indices from indices divisible by a prime. The
positivity discussion in [119] prompted careful distinction between
weighted and unweighted row polynomials. These are research prompts;
none of their unverified claims is used as a mathematical premise.
The primary references supplying the actual prerequisites are stated
at their points of use.

Two auxiliary questions guide the work. Can a critical free-Poisson
density have low positive values arbitrarily close to its necessarily
unbounded singularity? Can a nonconstant homogeneous rational weight
preserve an all-level cubic binomial congruence? The first is answered
by an explicit lacunary jump measure, and the second by the theorem in
Section~\ref{sec:general-weight}, with all hypotheses verified for
A176285. The applications keep the original named problems unchanged.
They do not establish irrationality of a new constant, nor are the
chosen conjectures claimed to be more important than every result in
the source corpus.


% Research note for internal mathematical audit, November 15, 2026.
% Original source: Hasebe--Sakuma, arXiv:1508.01285v2, Conjecture 3.7.
% NOT a claim of established historical priority.
\section{Critical free compound Poisson densities}\label{sec:free}

Let $\nu$ be any Borel probability measure on $\mathbb R$ with
$\nu(\{0\})=0$, and let $\mu=\pi\boxtimes\nu$ be the free compound
Poisson law of rate one.  Hasebe--Sakuma \cite[Conjecture 3.7]{HS}, asks whether
the density of $\mu$ is discontinuous at zero, and more strongly whether
it tends to infinity there.  There is no positive-support hypothesis on
$\nu$ in this conjecture; their equation (2.14) defines the law for every
probability measure on $\mathbb R$.

Write $G_\lambda(z)=\int(z-t)^{-1}\lambda(dt)$ and
$F_\lambda=1/G_\lambda$.  The analytic continuation of $F_\mu^{-1}$ is
\[
 H(z)=z+\int\frac{zt}{z-t}\nu(dt)=z^2G_\nu(z).
\]
For $u\in\mathbb R$, $v>0$, put
\[
 g(u,v)=\int\frac{t^2}{(u-t)^2+v^2}\nu(dt),\qquad
 v(u)=\inf\{v>0:g(u,v)<1\}.
\]
Then $\operatorname{Im}H(u+iv)=v(1-g(u,v))$.
We use the established inverse-domain theorem and boundary-density
formula of Huang, recorded as (2.16)--(2.19) and Theorem 2.4 in
Hasebe--Sakuma \cite{HS}: $v$ is continuous,
\[
 F_\mu(\mathbb C^+)=\Omega=\{u+iv:v>v(u)\},\qquad
 \psi(u)=H(u+iv(u))
\]
is a homeomorphism of $\mathbb R$, and the canonical density
is
\[
 \rho(\psi(u))=\frac{v(u)}{\pi(u^2+v(u)^2)}\quad(u\ne0).
\]
The density is analytic and positive on $\psi(V)$, where
$V=\{u:v(u)>0\}$.  Theorem 3.4(2) in Hasebe--Sakuma gives absolute
continuity of $\mu$.  Notice that $v(0)=0$ and $H(i0)=0$, so $\psi(0)=0$.

\begin{theorem}\label{thm:critical-unboundedness}
For every such $\nu$ and every $\delta>0$,
\[
 \mathop{\rm ess\,sup}_{|x|<\delta}\rho(x)=+\infty.
\]
In particular no representative of the density is continuous at zero
with a finite value.
\end{theorem}
\begin{proof}
For $y>0$ one has
\[
 \operatorname{Im}H(iy)=y^3\int\frac{\nu(dt)}{t^2+y^2}>0,
\]
and $iy\in\Omega$.  Consequently $z_y:=H(iy)\in\mathbb C^+$ and
$G_\mu(z_y)=1/(iy)$, whence $-\operatorname{Im}G_\mu(z_y)=1/y$.
Also $|H(iy)|=y^2|G_\nu(iy)|\le y$, so $z_y\to0$.
Suppose that $\rho\le M$ almost everywhere on $(-\delta,\delta)$.
For $z=a+ib\in\mathbb C^+$ with $|a|<\delta/2$, the Poisson integral
then satisfies
\[
 -\operatorname{Im}G_\mu(z)
 =\int\frac{b}{(a-x)^2+b^2}\rho(x)\,dx
 \le \pi M+\frac{4b}{\delta^2}.
\]
Apply this at $z_y$ and let $y\downarrow0$ to obtain a contradiction.
\end{proof}

The stronger asserted limit is false even if the limit is restricted to
the open set on which the canonical density is positive.
The construction below uses the established idea of locating gaps through
the boundary integral.  Huang \cite[Proposition 4.13]{Huang15} and Hasebe--Sakuma \cite[Example 5.3]{HS} use this idea with separated atoms
escaping to infinity.  Here the atoms accumulate at zero and the critical
total mass is exactly one; the quantitative outer-mass/inner-mass estimate
below is needed to produce gaps at the singularity itself.

\begin{theorem}\label{thm:critical-oscillation}
Set $a_n=2^{-2^n}$ and $w_n=2^{-n}$ for $n\ge1$, and define
\[
 \nu=\sum_{n\ge1}w_n\delta_{a_n},\qquad
 \mu=\pi\boxtimes\nu.
\]
Then $\mu$ is absolutely continuous and its density is locally
essentially unbounded at zero.  Nevertheless there are open gaps in
$\operatorname{supp}\mu$ accumulating at zero, and there is a sequence
$x_n\to0$ in the open set where its canonical density is positive with
$\rho(x_n)\to0$.
\end{theorem}
\begin{proof}
The measure $\nu$ is a compactly supported probability on $[0,\infty)$
with full mass on $(0,\infty)$ and no atom at zero.  For $u$ away from its atoms,
\[
 g(u,0)=\sum_{k\ge1}w_k\frac{a_k^2}{(u-a_k)^2}.
\]
For $n\ge5$ put
\[
 A=a_{n-1},\quad W=\sum_{k\ge n}w_k=2^{-(n-1)},\quad
 u_n=A^{3/2}=\sqrt{a_{n-1}a_n}.
\]
We have $A\le2^{-16}$ and $\sqrt A\le W/8$.  The latter inequality is
equivalent to $2^{n-2}\ge n+2$, valid at $n=5$ and preserved by
induction.  If $k<n$, then $a_k\ge A$ and
$r=u_n/a_k\le\sqrt A\le1/256$.  Since for $0\le r\le1/16$,
\[
 \frac{1}{(1-r)^2}-1
 =\frac{r(2-r)}{(1-r)^2}\le4r,
\]
the sum of these terms is at most $1-W+4\sqrt A$.
If $k\ge n$, then $a_k\le A^2$, and $s=a_k/u_n\le\sqrt A$.
Since $s^2/(1-s)^2\le2s^2\le2A$ in this range, their sum is at most
$2AW$.  Thus
\[
 g(u_n,0)\le1-W+4\sqrt A+2AW<1,
\]
because $4\sqrt A\le W/2$ and $2A<1/2$.

The point $u_n$ lies strictly between $a_n$ and $a_{n-1}$ and is separated
from the support of $\nu$, so $g(u,0)$ is continuous near $u_n$.
Hence $g(u,0)<1$ on an open interval $I_n$ containing $u_n$, and
$v(u)=0$ there.  Its image $\psi(I_n)$ is an open interval on which
$\rho=0$, hence a gap in the support.  Since $u_n\to0$ and
$\psi(0)=0$, these gaps accumulate at zero.

For completeness, the gaps do not merely exploit approach from outside
the support.  At $u=a_n$, the integral $g(a_n,0)$ is infinite, so
$v(a_n)>0$.  Let $J_n$ be the connected component of $V$ containing
$a_n$, and let $b_n$ be its right endpoint.  The gap around $u_n$ shows
that $a_n<b_n<u_n$.  Continuity of $v$ gives $v(b_n)=0$, and, as
$u\uparrow b_n$ within $J_n$,
\[
 \rho(\psi(u))=\frac{v(u)}{\pi(u^2+v(u)^2)}\longrightarrow0,
\]
since $b_n>0$.  Choose $q_n\in J_n\cap(a_n,b_n)$ sufficiently close to
$b_n$ that $0<\rho(\psi(q_n))<1/n$.  The points
$x_n=\psi(q_n)$ belong to the positive analytic-density set and tend to
zero because $0<q_n<u_n\to0$.  This proves the claim.
\end{proof}

Combining the two theorems gives the precise oscillatory conclusion
\[
 \liminf_{\substack{x\to0\\\rho(x)>0}}\rho(x)=0,
 \qquad
 \limsup_{\substack{x\to0\\\rho(x)>0}}\rho(x)=+\infty.
\]
Indeed the first equality was constructed above, whereas the second
follows from local essential unboundedness and continuity away from zero.



\section{Bala's cubic supercongruence}\label{sec:bala}\subsection{The unchanged sequence and conjecture}
The original sequence A176285 is defined by $a(1)=1$ and
\begin{equation}\label{original}
4(2n+1)^2a(n+1)+n^2a(n)
=(205n^2+160n+32)\binom{2n-1}{n}^{\!3}\quad(n\ge1).
\end{equation}
The OEIS entry \cite{OEIS}, in its formula block attributed to Peter Bala, dated 11 July 2024, records
\begin{equation}\label{def}
 a(N)=\sum_{k=0}^{N-1}
 \binom{N-1}{k}^{\!2}\binom{N+k-1}{k}
 \binom{2N+k-1}{N-1}\qquad(N\ge1).
\end{equation}
It conjectures $a(np^r)\equiv a(np^{r-1})\pmod{p^{3r}}$ for every
prime $p\ge5$ and all $n,r\ge1$.

\begin{theorem}\label{main}
For every prime $p\ge5$ and every positive integer $N$ divisible by $p$,
\[
 a(N)\equiv a(N/p)\pmod{p^{3\vp(N)}}.
\]
Consequently Bala's conjecture holds with all its original quantifiers.
\end{theorem}

\subsection{The two kernels}
Put
\begin{align*}
 U(N,k)&=\binom Nk^{\!2}\binom{N+k}{k}\binom{2N+k}{N},
 &0\le k\le N,\\
 w(N,k)&=\frac{(N-k)^2}{(N+k)(2N+k)},
 &T(N,k)&=w(N,k)U(N,k).
\end{align*}
Thus $T(N,N)=0$. For $0\le k<N$, factorial cancellation gives
\begin{align}
 U(N,k)&=\frac{(2N+k)!}{k!^3(N-k)!^2},\label{Ufact}\\
 T(N,k)&=\frac{(2N+k-1)!}{k!^3(N-k-1)!^2(N+k)}\notag\\
 &=\binom{N-1}{k}^{\!2}\binom{N+k-1}{k}
   \binom{2N+k-1}{N-1}.\label{Tintegral}
\end{align}
In particular every $T(N,k)$ is an integer and
$a(N)=\sum_{k=0}^NT(N,k)$. The rational weight satisfies
\begin{equation}\label{weight}
 w(pM,pj)=w(M,j),\qquad
 w(N,k)-1=-\frac{N(N+5k)}{(N+k)(2N+k)}.
\end{equation}

\begin{lemma}[Pfaff--Saalsch\"utz evaluation]\label{total}
For every nonnegative integer $N$,
\[
 \sum_{k=0}^N U(N,k)=\binom{2N}{N}^{\!3}.
\]
\end{lemma}
\begin{proof}
After dividing by $\binom{2N}{N}$, the left side is
\[
 {}_3F_2\!\left(\begin{matrix}-N,-N,2N+1\\1,1\end{matrix};1\right).
\]
The terminating Pfaff--Saalsch\"utz formula gives
\[
 \frac{(N+1)_N(-2N)_N}{(1)_N(-N)_N}
 =\binom{2N}{N}^{\!2}.
\]
All denominators in this finite identity are nonzero, including at
$N=0$ by the empty-product convention. This is the classical terminating
Pfaff--Saalsch\"utz identity \cite[Eq.~16.4.3]{DLMF} with parameters
$a=-N$, $b=2N+1$, $c=1$, and termination index $N$.
\end{proof}

\begin{lemma}[Exact identification with the original sequence]\label{identify}
The sum \eqref{def} is the sequence defined by \eqref{original} and $a(1)=1$.
\end{lemma}
\begin{proof}
Set $b(N)=\sum_{k=0}^NT(N,k)$, so that $b(1)=1$. For $n\ge1$ put
\[
P(n)=205n^2+160n+32,\qquad
R(n,k)=\frac{k^3(55n+23k)}{8(n+k)(2n+k)}.
\]
The following exact telescoping identity holds for $0\le k\le n$:
\begin{align}
&4(2n+1)^2T(n+1,k)+n^2T(n,k)-\frac{P(n)}8U(n,k)\notag\\
&\hspace{15mm}=R(n,k+1)U(n,k+1)-R(n,k)U(n,k),\label{certificate}
\end{align}
where $U(n,n+1)=0$. For a direct verification, divide by $U(n,k)$ and use
\[
\frac{T(n+1,k)}{U(n,k)}=\frac{2n+k+1}{n+k+1},\qquad
\frac{U(n,k+1)}{U(n,k)}=\frac{(2n+k+1)(n-k)^2}{(k+1)^3}.
\]
The resulting rational identity follows after multiplication by
$8(n+k)(2n+k)(n+k+1)$ and expansion; these denominator factors are
nonzero for the stated range. Summing \eqref{certificate} telescopes,
since $R(n,0)=0$ and $U(n,n+1)=0$. Also $T(n+1,n+1)=0$, so the first sum
on the left equals $b(n+1)$. Lemma~\ref{total} now gives
\[
4(2n+1)^2b(n+1)+n^2b(n)
=\frac{P(n)}8\binom{2n}{n}^{\!3}
=P(n)\binom{2n-1}{n}^{\!3}.
\]
The first-order recurrence and initial value uniquely determine the
sequence, proving $b(N)=a(N)$ for all $N\ge1$.
\end{proof}

We use the classical Jacobsthal congruence in the form recorded as
Lemma 5.1 of Straub \cite{Straub}: for $p\ge5$ and integers $A>B>0$,
\begin{equation}\label{jac}
 \frac{\binom{pA}{pB}}{\binom AB}
 \in 1+p^{3+\vp(AB(A-B))}\mathbb Z_p.
\end{equation}
The cases $B=0,A$ have ratio one. Ratios here are rational numbers
viewed in $\mathbb Q_p$; the assertion includes their being units.

\subsection{The divisible indices}
\begin{lemma}\label{divide}
Let $p\ge5$, $p\mid N$, and $\alpha=\vp(N)$. For every $k$ with
$0\le k\le N$ and $p\mid k$,
\begin{align*}
 U(N,k)&\equiv U(N/p,k/p)\pmod{p^{3\alpha}},\\
 T(N,k)&\equiv T(N/p,k/p)\pmod{p^{3\alpha}}.
\end{align*}
\end{lemma}
\begin{proof}
First suppose $0<k<N$, and put $s=\vp(k)\ge1$.
The four binomial factors of $U$, with the first repeated twice,
give Jacobsthal quotient precisions respectively
\[
 \vp(Nk(N-k)),\quad
 \vp((N+k)Nk),\quad
 \vp((2N+k)N(N+k)).
\]
If $s<\alpha$, all three quantities equal $\alpha+2s$.
It follows that
\[
 U(N,k)=U(N/p,k/p)R,\quad
 R\in1+p^{\alpha+2s}\mathbb Z_p.
\]
Because the weight is homogeneous, the same ratio $R$ satisfies
$T(N,k)=T(N/p,k/p)R$.
The identity $\binom AB=(A/B)\binom{A-1}{B-1}$ shows
\[
 \vp\binom{N/p}{k/p}\ge\alpha-s.
\]
The weight $w(N/p,k/p)$ is a $p$-adic unit: the valuations of
$N/p-k/p$, $N/p+k/p$ and $2N/p+k/p$ are all $s-1$.
Consequently both $U(N/p,k/p)$ and $T(N/p,k/p)$ have valuation
at least $2(\alpha-s)$. Multiplying by $R-1$ proves the required
valuation $2(\alpha-s)+\alpha+2s=3\alpha$.

If $s\ge\alpha$, each of the three displayed precisions is at
least $3\alpha$. Both lower-scale terms are integers by their
binomial formulas. Multiplication by $R-1\in p^{3\alpha}\mathbb Z_p$
therefore proves both congruences, whether or not the rational
weight by itself has a denominator divisible by $p$.

At $k=0$, $U(N,0)=\binom{2N}{N}$ and $T(N,0)=\tfrac12\binom{2N}{N}$.
Jacobsthal and the fact that 2 is a unit give the result.
At $k=N$, the $T$ terms vanish, while
$U(N,N)=\binom{2N}{N}\binom{3N}{N}$.
Both factors have Jacobsthal precision at least $3\alpha$.
\end{proof}

\subsection{The primitive indices and the proof}
\begin{lemma}\label{primitive}
Under the hypotheses of Lemma~\ref{divide},
\[
 \sum_{\substack{0\le k\le N\\p\nmid k}}U(N,k)
 \equiv
 \sum_{\substack{0\le k\le N\\p\nmid k}}T(N,k)
 \equiv0\pmod{p^{3\alpha}}.
\]
\end{lemma}
\begin{proof}
Lemma~\ref{total} and Jacobsthal imply
\[
 \sum_{k=0}^NU(N,k)
 \equiv\sum_{j=0}^{N/p}U(N/p,j)\pmod{p^{3\alpha}}.
\]
Subtract the first congruence of Lemma~\ref{divide}, summed over
all divisible indices. The remaining sum of primitive terms is zero.

For each $p\nmid k$, the inequality
$\vp\binom Nk\ge\alpha$ gives $\vp U(N,k)\ge2\alpha$.
Both factors in the denominator of \eqref{weight} are units, so
$\vp(w(N,k)-1)\ge\alpha$. Thus $T(N,k)-U(N,k)$ is divisible
by $p^{3\alpha}$ termwise. Summing gives the second assertion.
\end{proof}

\begin{proof}[Proof of Theorem~\ref{main}]
Split $a(N)=\sum T(N,k)$ into divisible and primitive indices.
The former sum is congruent to $a(N/p)$ by Lemma~\ref{divide};
the latter is zero by Lemma~\ref{primitive}. For $N=np^r$ one has
$\alpha=r+\vp(n)\ge r$, so the original modulus $p^{3r}$ follows.
\end{proof}


\subsection{A general homogeneous-weight transfer}\label{sec:general-weight}
The proof above answers a broader auxiliary problem. The criterion in
the next theorem is proved in full, and its application to A176285 has
already been verified; it is not used as a substitute for that application.

\begin{theorem}\label{thm:general-weight}
Fix a prime $p\ge5$. Let $A(X,Y)$ and $B(X,Y)$ be homogeneous
polynomials of the same degree with integer coefficients, and suppose
$p\nmid B(0,1)$. Put $W=A/B$. Assume that $B(N,k)\ne0$ and
\[
 V(N,k):=W(N,k)U(N,k)\in\mathbb Z_p
 \quad(N\ge1,\ 0\le k\le N).
\]
Then $S_W(N)=\sum_{k=0}^N V(N,k)$ satisfies, for every $p\mid N$,
\[
 S_W(N)\equiv S_W(N/p)\pmod{p^{3v_p(N)}}.
\]
\end{theorem}
\begin{proof}
Write $\alpha=v_p(N)$ and $c=A(0,1)/B(0,1)\in\mathbb Z_p$.
When $p\nmid k$, the denominator $B(N,k)$ is a unit. Homogeneity
and reduction modulo $N$ give
\[
 W(N,k)-c\in p^\alpha\mathbb Z_p.
\]
Indeed $A(X,Y)B(0,1)-A(0,1)B(X,Y)$ is divisible by $X$.
Since $v_pU(N,k)\ge2\alpha$, the primitive terms obey
$V(N,k)\equiv cU(N,k)\pmod{p^{3\alpha}}$. Their sum vanishes
by Lemma~\ref{primitive}.

For $0<k<N$ divisible by $p$, put $s=v_p(k)$. If $s<\alpha$,
homogeneity gives
\[
 W(N/p,k/p)=W(N/p^s,k/p^s)\in\mathbb Z_p,
\]
because the first argument is divisible by $p$, the second is a unit,
and $B(0,1)$ is a unit. Thus
$v_p V(N/p,k/p)\ge2(\alpha-s)$. The binomial ratio from the
proof of Lemma~\ref{divide} is in
$1+p^{\alpha+2s}\mathbb Z_p$, so the difference of the two
$V$ terms has valuation at least $3\alpha$.
If $s\ge\alpha$, that ratio is in $1+p^{3\alpha}\mathbb Z_p$,
and the assumed integrality of the smaller-scale $V$ gives the same
conclusion. At $k=0,N$ the binomial ratios again have precision at
least $3\alpha$, and homogeneity and integrality complete the argument.
Adding the divisible terms and the vanishing primitive sum proves
the theorem.
\end{proof}

For A176285 take $A=(X-Y)^2$ and $B=(X+Y)(2X+Y)$. Here
$B(0,1)=1$, and the integral realization is exactly
\eqref{Tintegral}; no assumption of the transfer theorem remains open.
The gain $2\alpha+\alpha=3\alpha$ is the arithmetic content of the
weight removal. The mere identity $T=wU$ would not supply it.
The general idea of transferring binomial congruences is classical;
historical novelty of this precise weight class is not asserted.


\section{Proof audit and historical scope}\label{sec:audit}
All universal mathematical assertions in the two main sections have
proofs above using established results. The remaining qualification is
historical: a finite literature search cannot certify that no equivalent
proof or theorem has ever appeared.

\begin{center}\small
\begin{tabular}{@{}p{.08\textwidth}p{.84\textwidth}@{}}
\toprule
Gate & Assessment\\
\midrule
G1 & The named source statements predate this work: Hasebe--Sakuma
Conjecture 3.7 and Bala's dated OEIS conjecture.\\[4pt]
G2 & The free-Poisson law is unchanged. For Bala's sequence,
Lemma~\ref{identify} verifies equality with its original recurrence.
The stronger density statement is refuted even when approach is
restricted to positive-density points.\\[4pt]
G3 & Original sources label the statements conjectural. The inspected
current sources did not reveal a prior full resolution. A global claim
that they remained open immediately before this work is not certified.\\[4pt]
G4 & The main arguments are unconditional and cover every stated
parameter. Finite computations are checks, not premises of the proofs.\\[4pt]
G5 & The substantive estimates are $g(u_n,0)<1$ at the critical
singularity and $v_p((w-1)U)\ge3v_p(N)$ on primitive indices.
The weak discontinuity assertion alone is a short consequence of known
inverse theory. The Jacobi--Stirling appendix is an elementary finite
obstruction, not a new general method. Historical originality of the
methods is not certified.\\[4pt]
G6 & Not certified. Closely related earlier methods are credited,
and absence of a search hit is not treated as proof of first priority.\\
\bottomrule
\end{tabular}
\end{center}

\subsection{Primary-source comparisons}
For the density problem we inspected the original Conjecture 3.7,
its Theorem 2.4 and Theorem 3.4, and the later regularity treatments
of Huang--Wang \cite{HW} and Bercovici--Wang--Zhong \cite{BWZ}.
Those works provide the boundary homeomorphism and density formulas;
we do not claim to have invented them. The separated-atom construction
of support gaps also has the precedents explicitly cited in
Section~\ref{sec:free}. The specific estimate here places infinitely
many gaps at the critical origin, while producing arbitrarily small
positive density values inside the support.

For the arithmetic problem, Straub's multivariate congruence
\cite{Straub} does not directly cover the shifted ray
$(N-1,N-1,N-1,2N-1)$: division by $p$ does not preserve that shift.
The weight comparison and the exact recurrence certificate avoid this
unjustified application. The classical Pfaff--Saalsch\"utz evaluation
and the Jacobsthal estimate are explicitly attributed.
The existing proof of integrality for Sun's original sequence and
prime-level weighted congruences \cite{HP} are not claimed as new.
The result proved here is the particular all-level adjacent-scale
congruence stated by Bala.

The literature check on 15 November 2026 included exact target numbers,
authors, formula variants, and equivalent terminology, through web and
scholarly search, followed by the primary texts just identified.
The A176285 search also checked the public Bala-congruence portfolio
\cite{Bajaj} and Easwar's binomial-identity paper \cite{Easwar}.
No identical A176285 all-level theorem was located in those inspected
materials. This is a reproducible scope statement, not an exhaustive
priority determination.

\subsection{Independent mathematical checks}
The free-probability proof was reviewed separately for the inverse
branch, the domain of the jump law, the local Poisson bound, the
infinite lacunary estimate, and the support-interior interpretation of
the limit. The arithmetic proof was separately checked for every
valuation stratum, including a rational weight with nonunit denominator,
both endpoints, and extra factors of $p$ in the multiplier $n$.
The telescoping certificate and negative Hankel determinant were also
expanded independently in exact arithmetic. None of these checks is
represented as external peer review or a proof-assistant certification.

We excluded a rederivation of Written on the Wall II 434a--b after
finding the same result and proof in corpus [065]. This illustrates
why a correct argument can fail the novelty requirement.

\appendix
% Self-contained mathematical fragment; requires amsmath, amssymb, hyperref.
% All variables below follow Zhu's original notation in equation (5.15).
\section{An elementary counterexample to the unweighted Jacobi--Stirling conjecture}
\label{sec:jacobi-stirling-counterexample}

This section gives a finite exact counterexample to a previously stated
conjecture.  It is an elementary certificate, not a claim of a new general
method of total positivity or of established historical priority.

Bao-Xuan Zhu defines the generalized Jacobi--Stirling numbers by
\begin{equation}\label{eq:js-original-recurrence}
 \mathbf{JS}_{n,k}=\mathbf{JS}_{n-1,k-1}
 +(xk^2+yk)\mathbf{JS}_{n-1,k},
 \qquad \mathbf{JS}_{0,0}=1,
\end{equation}
with entries zero outside $0\leq k\leq n$.  Immediately after equation
(5.15) in Section 5.2 of his paper, Zhu records his 2018 conjecture that
its unweighted row-generating polynomials are Hankel-totally positive in
$\mathbb Z[x,y,q]$ with the coefficientwise order.\footnote{B.-X. Zhu,
\emph{Coefficientwise Hankel-total positivity of the row-generating
polynomials for the output matrices of certain production matrices},
Advances in Mathematics 447 (2024), 109695, Section 5.2, p.~26;
\url{https://doi.org/10.1016/j.aim.2024.109695}.
The inspected primary preprint is
\url{https://arxiv.org/html/2202.03793v2}, equation (5.15) and the
following paragraph (version 2, 22 April 2024).}
Thus the precise target is
\[
 J_n(q)=\sum_{k=0}^n\mathbf{JS}_{n,k}q^k,
 \qquad H=(J_{i+j}(q))_{i,j\geq0}.
\]
Coefficientwise Hankel total positivity would require every minor of
$H$ to have nonnegative coefficients in the original three variables.

Set $c=x+y$, $d=4x+2y$, and $e=9x+3y$.  The original recurrence gives
\begin{align*}
 J_0&=1, &J_1&=q,\\
 J_2&=q^2+cq,\\
 J_3&=q^3+(c+d)q^2+c^2q,\\
 J_4&=q^4+(c+d+e)q^3+(c^2+cd+d^2)q^2+c^3q.
\end{align*}
Direct determinant expansion now yields the polynomial identity
\begin{align}
 \det\!\begin{pmatrix}J_0&J_1&J_2\\J_1&J_2&J_3\\J_2&J_3&J_4\end{pmatrix}
 &=q^3\{cd(d-c)+q(cd+ce-c^2-d^2)\}\notag\\
 &=2q^3\bigl(6x^3+11x^2y+6xy^2+y^3-2x^2q\bigr).
 \label{eq:js-negative-minor}
\end{align}
The coefficient of $x^2q^4$ is $-4$.  Consequently the unweighted
conjecture is false.  This is an exact finite identity and requires no
numerical approximation, asymptotic hypothesis, or unverified limiting
argument.

There is also a numerical certificate at the permitted specialization
$(x,y,q)=(1,0,4)$:
\[
 (J_0,J_1,J_2,J_3,J_4)=(1,4,20,148,1492),\qquad
 \det\!\begin{pmatrix}
 1&4&20\\4&20&148\\20&148&1492
 \end{pmatrix}=-256.
\]
The failure is not confined to a boundary of the nonnegative parameter
region: at $(x,y,q)=(1,1,13)$ the same determinant is $-8788$.

\paragraph{Distinction from the factorial-weighted theorem.}
Zhu's Proposition 5.7 proves Hankel total positivity for
$\widetilde J_n(q)=\sum_k\mathbf{JS}_{n,k}k!q^k$ (and its reversal).
That factorial-weighted statement is different from the conjecture just
refuted.  No contradiction with Proposition 5.7 is asserted here, and
no weighting or normalization is inserted into the target.

\paragraph{A structural obstruction.}
For the more general recurrence
$T_{n,k}=T_{n-1,k-1}+\lambda_kT_{n-1,k}$ with $\lambda_0=0$, the same
calculation shows that the coefficient of $q^4$ in the leading
$3\times3$ row-polynomial Hankel minor is
\[
 \lambda_1\lambda_2+\lambda_1\lambda_3
 -\lambda_1^2-\lambda_2^2.
\]
For $\lambda_k=k^p$ with any real $p>1$, this is
$2^p+3^p-1-4^p<0$: indeed strict convexity of $t^p$ implies
$4^p-3^p>2^p-1$.  Hence every such superlinear-power spectrum fails
coefficientwise Hankel total positivity of its unweighted row sequence.
This observation explains a broad obstruction, but it is still an
algebraic consequence of the displayed finite minor.


\begin{thebibliography}{99}
\bibitem{HS} T. Hasebe and N. Sakuma, \emph{Unimodality for free L\'evy
processes}, Ann. Inst. H. Poincar\'e Probab. Statist. 53 (2017), 916--936;
Conjecture 3.7. \url{https://arxiv.org/abs/1508.01285}.
\bibitem{Huang15} H.-W. Huang, \emph{Supports of measures in a free
additive convolution semigroup}, Int. Math. Res. Not. (2015),
4269--4292; Proposition 4.13. \url{https://arxiv.org/abs/1205.5542}.
\bibitem{HW} H.-W. Huang and J.-C. Wang, \emph{Regularity results for
free L\'evy processes}, Adv. Math. 402 (2022), 108323.
\url{https://doi.org/10.1016/j.aim.2022.108323}.
\bibitem{BWZ} H. Bercovici, J.-C. Wang and P. Zhong,
\emph{Superconvergence to freely infinitely divisible distributions},
Pacific J. Math. (2018). \url{https://arxiv.org/abs/1510.03393}.
\bibitem{OEIS} OEIS Foundation, \emph{A176285}, sequence introduced
by Z.-W. Sun; formula and all-level conjecture by P. Bala,
11 July 2024. \url{https://oeis.org/A176285}.
\bibitem{DLMF} NIST Digital Library of Mathematical Functions,
Eq.~16.4.3, Pfaff--Saalsch\"utz balanced sum.
\url{https://dlmf.nist.gov/16.4#E3}.
\bibitem{Straub} A. Straub, \emph{Multivariate Ap\'ery numbers and
supercongruences of rational functions}, Algebra Number Theory
8 (2014), 1985--2008; Lemma 5.1.
\url{https://arxiv.org/abs/1401.0854}.
\bibitem{HP} Kh. Hessami Pilehrood and T. Hessami Pilehrood,
\emph{Bivariate identities for values of the Hurwitz zeta function and
supercongruences}, 2011 preprint, Section 5.
\url{https://arxiv.org/abs/1104.3659}.
\bibitem{Bajaj} R. Bajaj, \emph{A183068 supercongruence and related
results}, public research repository, consulted 15 November 2026.
\url{https://github.com/rbajaj5/a183068-supercongruence}.
\bibitem{Easwar} R. Easwar, \emph{An application of Jacobi's residue
formula and proofs of some other conjectures on binomial coefficients},
2026. \url{https://arxiv.org/abs/2608.22053}.
\bibitem{VMRC} QEEDDD, \emph{Verifiable Mathematical Reasoning Corpus},
snapshot \texttt{38348353416f1f7f}, consulted 15 November 2026.
\url{https://github.com/QEEDDD/Verifiable-Mathematical-Reasoning-Corpus}.
\end{thebibliography}
\end{document}
